\documentclass[a4paper,11pt,fleqn]{scrartcl}%fleqn pour aligner les equations à gauche
\usepackage{../../Styles/Cours}


\newcommand{\chnum}{Chapitre 05}
\newcommand{\chtitle}{Détermination d'une quantité de matière par titrage}
\newcommand{\dtype}{Cours}

\begin{document}
	\maketitle
\section{Rappels}
Une solution est un mélange liquide homogène de deux ou plusieurs substances. La substance majoritaire d'une solution est appelée le solvant, les autres substances dispersées dans le solvant sont appelées les solutés.\\
La quantité de matière $n$ d'un soluté de concentration $C$ dans une solution de volume $V$ est donnée par la relation : $n= C \cdot V$
\begin{exercice}
    On dissout une masse $m=\SI{10,0}{g}$ de chlorure de fer (III) dans de l'eau distillée afin d'obtenir un volume de solution $V = \SI{250}{mL}$ contenant du chlorure de fern (\ce{Fe^3+ ; Cl^-}).
    \begin{enumerate}
        \item Quel nom donne-t-on à une telle opération ?
        \item Déterminer la concentration en masse $\gamma$ de cette solution.
        \item Déterminer la concentration molaire $C$ de cette solution sachant que la masse molaire du chlorure de fer (III) est $M(\ce{FeCl3}) = \SI{162,5}{g.mol^{-1}}$.
    \end{enumerate}
\end{exercice}
\section{Principe d'un titrage}
Un titrahe, ou dosage, consiste à déterminer la concentration en quantité de matière $C$ ou la concentration en masse $\gamma$ d'un soluté dissout dans une solution.

\begin{exercice}
    Soient A et B deux espèces chimiques pouvant réagir selon l'équation bilan : \ce{A + 2B -> AB2}.\\
    On précise que l'espèce B est jaune en solution aqueuse alors que les espèces \ce{A} et \ce{AB2} sont inclores.\\
    On place dans un bécher un volume $V_A = \SI{100}{mL}$ d'une solution aqueuse $S_A$ contennat une quantité de matière inconnue $n_A$ de l'espèce chimique \ce{A}.\\
    On plmace ensuite, au-dessus du bécher, une burette graduée contenant une solution $S_B$ jaune de concentration en quantité de matière $C_B = \SI{0,0200}{mol.L^{-1}}$ de l'espèce chimique \ce{B}. Pour finir, on allume k'agitateur magnétique pour faire tourner un barreau aimanté placé dans le bécher.
    \begin{enumerate}
        \item Quel est le rôle du barreau aimanté ?
        \item On verse \SI{1,0}{mL} de la solution $S_B$ dans le bécher. Le volume introduit se décolore instantanémant en se mélangeant à la solution $S_A$. Expliquer cette observation.
        \item Combien de moles de \ce{B} faut-il pour consommer une mole de \ce{A} ?
        \item En déduire une relation  mathématique entre la quantité de matuère $n_B$ à introduire dans le bécher avec la burette pour faire réagir intégralement la quantité de matière $n_A$ initialement présente dans le bécher.
        \item Après avoir versé \SI{5,0}{mL} de la solution $S_B$, le contenu du bécher reste toujours incolore. Mais si l'on rajoute alors une goutte supplémentaire de solution $S_B$, on observe dans le bécher une légère coloration jaune qui persiste. Expliquer cette observation.
        \item Déterminer alors la quantité de matière $n_B$ introduire dans le bécher pou observer cette coloration persistante.
        \item En déduire la quantité de matière $n_A$ initialement présente dans le bécher.
        \item Déterminer alors la concentration en quantité de matière $C_A$ en espèce \ce{A} de la solution $S_A$.
        \item Pourquoir peut-on qualifier ce type de titrage de solorimétrique ?
    \end{enumerate}
\end{exercice}

\end{document}